% infra-glue — Glue: counts for arbitrary norm n (beyond the thesis)
%
% Prose twin: proofs/infra-glue.md. Ground truth for the per-prime inputs is
% ERRATA.md (E4 corrected Thm 3.1.2, E5 corrected Cor 3.1.3). Nothing in this
% file is stated in the thesis itself: the thesis counts only prime-power norms.
%
% Standing notation (restated from \ref{notation}): $d < 0$, $d \equiv 0, 1
% \pmod 4$, $d = Df^2$ with $D$ fundamental and $f \geq 1$ the conductor;
% $K = \mathbb{Q}(\sqrt d)$; $\tau_d = (d + \sqrt d)/2$; $\mathcal{O}_d =
% \mathbb{Z}[\tau_d] = \mathbb{Z} + f\mathcal{O}_K$; $\mathcal{O}_K =
% \mathcal{O}_D$ the maximal order. The norm of a nonzero ideal is the index
% $N(\mathfrak{a}) = [\mathcal{O}_d : \mathfrak{a}]$. We write
% $r_d(n)$ (resp.\ $r_d^{\times}(n)$) for the number of ideals (resp.\
% invertible ideals) of index $n$ in $\mathcal{O}_d$, and $r_D(n)$ for the
% count in $\mathcal{O}_K$; $w = v_p(f)$; $v = v_p(d/4)$ means $v_p(d)$ for
% odd $p$ and $v_2(d) - 2$ for $p = 2$; $\varepsilon(m) = 1$ if $m$ is even,
% $0$ if $m$ is odd (thesis convention); split/inert/ramified always refers to
% the behavior of $p$ in $\mathcal{O}_K$, i.e.\ to the Kronecker symbol
% $\chi_D(p)$.

\begin{proposition}[Prime-power factorization of ideal counts]
  \label{infra-glue-primary}
  \lean{QuadraticOrder.idealCount_eq_prod_primeFactors,
        QuadraticOrder.invertibleIdealCount_eq_prod_primeFactors}
  \uses{notation, infra-index}
  Let $d$ be an imaginary quadratic discriminant, $\mathcal{O} = \mathcal{O}_d$,
  and $n \geq 1$. Sending an ideal to its family of primary components is a
  bijection
  \[
    \{\mathfrak{a} \subseteq \mathcal{O} : [\mathcal{O} : \mathfrak{a}] = n\}
    \;\longleftrightarrow\;
    \prod_{p \mid n}
    \{\mathfrak{b} \subseteq \mathcal{O} : [\mathcal{O} : \mathfrak{b}] = p^{v_p(n)}\},
  \]
  under which $\mathfrak{a}$ is invertible if and only if every component is
  invertible. Consequently
  \[
    r_d(n) = \prod_{p \mid n} r_d\bigl(p^{v_p(n)}\bigr),
    \qquad
    r_d^{\times}(n) = \prod_{p \mid n} r_d^{\times}\bigl(p^{v_p(n)}\bigr).
  \]
\end{proposition}

\begin{proof}
  \uses{notation, infra-index}
  We use the following index facts from the index layer, each elementary:
  (0.1) an ideal of finite index $N$ contains $N$ (Lagrange in
  $\mathcal{O}/\mathfrak{a}$);
  (0.2) if $\mathfrak{b} \subseteq \mathfrak{c}$ then
  $[\mathcal{O} : \mathfrak{c}] \mid [\mathcal{O} : \mathfrak{b}]$;
  (0.3) ideals of coprime finite indices are comaximal (by (0.2), the index of
  the sum divides both indices);
  (0.4) $[\mathcal{O} : \mathfrak{b}\mathfrak{c}]$ is finite and divides
  $N_{\mathfrak{b}}^2 N_{\mathfrak{c}}$, since
  $N_{\mathfrak{b}}\mathfrak{c} \subseteq \mathfrak{b}\mathfrak{c}$ by (0.1)
  and $[\mathcal{O} : N_{\mathfrak{b}}\mathfrak{c}] =
  N_{\mathfrak{b}}^2 N_{\mathfrak{c}}$;
  (0.5) comaximal ideals satisfy
  $\mathfrak{b} \cap \mathfrak{c} = \mathfrak{b}\mathfrak{c}$
  (write $1 = b + c$ and expand $x = xb + xc$);
  (0.6) the Chinese remainder theorem for finitely many pairwise comaximal
  ideals.

  \emph{Components.} Let $[\mathcal{O} : \mathfrak{a}] = n$ and
  $M := \mathcal{O}/\mathfrak{a}$, a finite $\mathcal{O}$-module of cardinality
  $n$. As a finite abelian group, $M = \bigoplus_{p \mid n} M_p$ with $M_p$ the
  $p$-primary component, of order $p^{v_p(n)}$. Each $M_p$ is an
  $\mathcal{O}$-submodule: multiplication by $y \in \mathcal{O}$ is an additive
  endomorphism, and $p^k x = 0$ implies $p^k(yx) = y(p^k x) = 0$. Hence the
  projections $\pi_p : M \to M_p$ are $\mathcal{O}$-linear, and
  \[
    \mathfrak{a}_p := \ker\bigl(\mathcal{O} \twoheadrightarrow M
      \xrightarrow{\pi_p} M_p\bigr)
  \]
  is an ideal containing $\mathfrak{a}$ with
  $\mathcal{O}/\mathfrak{a}_p \cong M_p$, so
  $[\mathcal{O} : \mathfrak{a}_p] = p^{v_p(n)}$; and
  $\bigcap_{p \mid n} \mathfrak{a}_p = \ker(\mathcal{O} \to \bigoplus_p M_p
  = M) = \mathfrak{a}$.

  \emph{Inverse map.} Given ideals $\mathfrak{b}_p$ with
  $[\mathcal{O} : \mathfrak{b}_p] = p^{v_p(n)}$, they are pairwise comaximal by
  (0.3), so by (0.6) the intersection
  $\mathfrak{b} := \bigcap_p \mathfrak{b}_p$ satisfies
  $\mathcal{O}/\mathfrak{b} \cong \prod_p \mathcal{O}/\mathfrak{b}_p$ and has
  index $n$. The $p$-primary component of the product is the factor
  $\mathcal{O}/\mathfrak{b}_p$ (a $p$-group; the other factors have order prime
  to $p$), so the component construction applied to $\mathfrak{b}$ returns
  exactly the $\mathfrak{b}_p$. The two constructions are therefore mutually
  inverse, and counting gives the product formula for $r_d$.

  \emph{Invertibility.} By (0.5) and induction,
  $\mathfrak{a} = \bigcap_p \mathfrak{a}_p = \prod_p \mathfrak{a}_p$. If each
  $\mathfrak{a}_p$ has a fractional inverse $\mathfrak{c}_p$, then
  $\mathfrak{a} \cdot \prod_p \mathfrak{c}_p = \mathcal{O}$. Conversely, if
  $\mathfrak{a}\mathfrak{c} = \mathcal{O}$, then for each $p$,
  $\mathfrak{a}_p \cdot \bigl(\mathfrak{c}\prod_{q \neq p}
  \mathfrak{a}_q\bigr) = \mathcal{O}$. So the bijection restricts to invertible
  ideals on both sides, giving the product formula for $r_d^{\times}$.
\end{proof}

\begin{lemma}[Conductor-avoiding extension and contraction]
  \label{infra-glue-avoid}
  \uses{notation, infra-index}
  Let $d = Df^2$ and let $N \geq 1$ with $\gcd(N, f) = 1$. Then:
  \begin{enumerate}
    \item $\mathcal{O}_d \cap N\mathcal{O}_K = N\mathcal{O}_d$, and the
      inclusion $\mathcal{O}_d \hookrightarrow \mathcal{O}_K$ induces a ring
      isomorphism
      $\mathcal{O}_d/N\mathcal{O}_d \xrightarrow{\ \sim\ }
       \mathcal{O}_K/N\mathcal{O}_K$;
    \item for every ideal $\mathfrak{c} \subseteq \mathcal{O}_d$ of index $N$:
      $\mathfrak{c}\mathcal{O}_K = \mathfrak{c} + N\mathcal{O}_K$,
      \quad
      $\mathfrak{c}\mathcal{O}_K \cap \mathcal{O}_d = \mathfrak{c}$,
      \quad and \quad
      $[\mathcal{O}_K : \mathfrak{c}\mathcal{O}_K] = N$.
  \end{enumerate}
\end{lemma}

\begin{proof}
  \uses{notation, infra-index}
  Fix $\alpha, \beta \in \mathbb{Z}$ with $1 = \alpha N + \beta f$, and recall
  $f\mathcal{O}_K \subseteq \mathcal{O}_d = \mathbb{Z} + f\mathcal{O}_K$.

  (1) If $x = Ny \in \mathcal{O}_d$ with $y \in \mathcal{O}_K$, then
  $y = \alpha Ny + \beta fy = \alpha x + \beta(fy) \in \mathcal{O}_d$ since
  $fy \in f\mathcal{O}_K \subseteq \mathcal{O}_d$; hence
  $x \in N\mathcal{O}_d$, proving
  $\mathcal{O}_d \cap N\mathcal{O}_K \subseteq N\mathcal{O}_d$ (the reverse
  inclusion is clear). The composite
  $\mathcal{O}_d \hookrightarrow \mathcal{O}_K \to
  \mathcal{O}_K/N\mathcal{O}_K$ is a ring homomorphism with kernel
  $\mathcal{O}_d \cap N\mathcal{O}_K = N\mathcal{O}_d$, so the induced map is
  injective; it is surjective because every $z \in \mathcal{O}_K$ satisfies
  $z = \alpha Nz + \beta fz \equiv \beta fz \pmod{N\mathcal{O}_K}$ with
  $\beta fz \in \mathcal{O}_d$.

  (2) First, $N \in \mathfrak{c}$ (Lagrange), so
  $N\mathcal{O}_K \subseteq \mathfrak{c}\mathcal{O}_K$ and
  $\mathfrak{c} + N\mathcal{O}_K \subseteq \mathfrak{c}\mathcal{O}_K$;
  conversely each generator $cz$ ($c \in \mathfrak{c}$, $z \in \mathcal{O}_K$)
  satisfies $cz = \alpha N(cz) + \beta c(fz) \in N\mathcal{O}_K + \mathfrak{c}$
  because $fz \in \mathcal{O}_d$ makes $c(fz) \in \mathfrak{c}$. This proves
  $\mathfrak{c}\mathcal{O}_K = \mathfrak{c} + N\mathcal{O}_K$.
  Next, if $x \in \mathfrak{c}\mathcal{O}_K \cap \mathcal{O}_d$, write
  $x = c + Ny$ with $c \in \mathfrak{c}$, $y \in \mathcal{O}_K$; then
  $Ny = x - c \in \mathcal{O}_d \cap N\mathcal{O}_K = N\mathcal{O}_d
  \subseteq \mathfrak{c}$ by (1) and Lagrange, so $x \in \mathfrak{c}$, proving
  $\mathfrak{c}\mathcal{O}_K \cap \mathcal{O}_d = \mathfrak{c}$.
  Finally, under the isomorphism of (1) the image of
  $\mathfrak{c}/N\mathcal{O}_d$ is
  $(\mathfrak{c} + N\mathcal{O}_K)/N\mathcal{O}_K =
  \mathfrak{c}\mathcal{O}_K/N\mathcal{O}_K$, so
  $[\mathcal{O}_K : \mathfrak{c}\mathcal{O}_K] =
   [\mathcal{O}_d/N\mathcal{O}_d : \mathfrak{c}/N\mathcal{O}_d] = N$.
\end{proof}

\begin{theorem}[Conductor-avoiding correspondence]
  \label{infra-glue-coprime}
  \lean{QuadraticOrder.idealCount_eq_of_coprime_conductor}
  \uses{notation, infra-index}
  Let $d = Df^2$ and let $n \geq 1$ with $\gcd(n, f) = 1$. Then extension
  $\mathfrak{a} \mapsto \mathfrak{a}\mathcal{O}_K$ and contraction
  $\mathfrak{A} \mapsto \mathfrak{A} \cap \mathcal{O}_d$ are mutually inverse,
  index-preserving bijections
  \[
    \{\mathfrak{a} \subseteq \mathcal{O}_d :
      [\mathcal{O}_d : \mathfrak{a}] = n\}
    \;\longleftrightarrow\;
    \{\mathfrak{A} \subseteq \mathcal{O}_K :
      [\mathcal{O}_K : \mathfrak{A}] = n\},
  \]
  and every ideal of $\mathcal{O}_d$ of index coprime to $f$ is invertible.
  In particular
  \[
    r_d(n) = r_d^{\times}(n) = r_D(n) = r_D^{\times}(n).
  \]
\end{theorem}

\begin{proof}
  \uses{infra-glue-avoid, notation, infra-index}
  \emph{Bijection.} Every ideal of index $n$ contains $n$ (Lagrange), so ideals
  of index $n$ in $\mathcal{O}_d$ (resp.\ $\mathcal{O}_K$) correspond to ideals
  of $\mathcal{O}_d/n\mathcal{O}_d$ (resp.\ $\mathcal{O}_K/n\mathcal{O}_K$)
  with quotient of cardinality $n$. The ring isomorphism of
  Lemma~\ref{infra-glue-avoid}(1) (with $N = n$) matches the two collections
  bijectively, preserving indices. By Lemma~\ref{infra-glue-avoid}(2) the
  forward map is $\mathfrak{a} \mapsto \mathfrak{a}\mathcal{O}_K$ (of index
  $n$), and the inverse is
  $\mathfrak{A} \mapsto \mathfrak{A} \cap \mathcal{O}_d$ (the preimage of
  $\mathfrak{A}/n\mathcal{O}_K$ in $\mathcal{O}_d$, using
  $n\mathcal{O}_K \subseteq \mathfrak{A}$); mutual inversion is
  $\mathfrak{c}\mathcal{O}_K \cap \mathcal{O}_d = \mathfrak{c}$ together with
  bijectivity of the mod-$n$ isomorphism.

  \emph{Invertibility.} Let $[\mathcal{O}_d : \mathfrak{a}] = n$ and
  $\mathfrak{A} := \mathfrak{a}\mathcal{O}_K$. In the maximal order,
  $\mathfrak{A}\bar{\mathfrak{A}} = N(\mathfrak{A})\mathcal{O}_K =
  n\mathcal{O}_K$ (standard; by multiplicativity this reduces to prime ideals,
  where it is checked in each splitting case — cf.\ Cox, Lemma~7.14).
  Conjugation is a ring automorphism of $\mathcal{O}_d$
  ($\bar\tau_d = d - \tau_d$), so $\bar{\mathfrak{a}}$ is an ideal of index
  $n$, and $[\mathcal{O}_d : \mathfrak{a}\bar{\mathfrak{a}}]$ is finite and
  divides $n^3$, hence is coprime to $f$. Applying
  Lemma~\ref{infra-glue-avoid}(2) to
  $\mathfrak{c} = \mathfrak{a}\bar{\mathfrak{a}}$, and using that extension is
  multiplicative and commutes with conjugation,
  \[
    \mathfrak{a}\bar{\mathfrak{a}}
    = \bigl(\mathfrak{a}\bar{\mathfrak{a}}\bigr)\mathcal{O}_K \cap \mathcal{O}_d
    = \mathfrak{A}\bar{\mathfrak{A}} \cap \mathcal{O}_d
    = n\mathcal{O}_K \cap \mathcal{O}_d
    = n\mathcal{O}_d,
  \]
  the last step by Lemma~\ref{infra-glue-avoid}(1). Hence
  $\mathfrak{a} \cdot \frac{1}{n}\bar{\mathfrak{a}} = \mathcal{O}_d$ and
  $\mathfrak{a}$ is invertible.

  \emph{Counts.} The bijection gives $r_d(n) = r_D(n)$; invertibility of every
  ideal in sight gives $r_d^{\times}(n) = r_d(n)$, and $\mathcal{O}_K$ is
  Dedekind, so $r_D^{\times}(n) = r_D(n)$.
\end{proof}

\begin{theorem}[Ideal counts for arbitrary norm: assembled closed form]
  \label{infra-glue}
  \uses{notation}
  Let $d = Df^2$ be an imaginary quadratic discriminant. For a prime $p$ and
  $m \geq 0$ define local factors as follows. If $p \nmid f$:
  \[
    t_p(m) = u_p(m) =
    \begin{cases}
      m + 1 & p \text{ split in } \mathcal{O}_K,\\
      \varepsilon(m) & p \text{ inert in } \mathcal{O}_K,\\
      1 & p \mid D.
    \end{cases}
  \]
  If $p \mid f$, with $w = v_p(f)$ and $v = v_p(d/4)$:
  \[
    t_p(m) =
    \begin{cases}
      \dfrac{p^{\lfloor m/2 \rfloor + 1} - 1}{p - 1} & m \leq v,\\[2ex]
      \dfrac{p^{w+1} - 1}{p - 1} & m > v,\ p \mid D,\\[2ex]
      \dfrac{p^{w+\varepsilon(m)} - 1}{p - 1} & m > v,\ p \text{ inert},\\[2ex]
      \dfrac{p^{w+\varepsilon(m)} - 1}{p - 1}
        + p^w\bigl(m + 1 - 2w - \varepsilon(m)\bigr) & m > v,\ p \text{ split},
    \end{cases}
  \]
  \[
    u_p(m) =
    \begin{cases}
      \varepsilon(m)\, p^{m/2}
        & m < v,\ \text{or } m = v = 2w - 2,\\
      p^w
        & (m > v \text{ or } m = v \geq 2w),\ p \mid D,\\
      \varepsilon(m)\bigl(p^w + p^{w-1}\bigr)
        & (m > v \text{ or } m = v \geq 2w),\ p \text{ inert},\\
      \bigl(p^w - p^{w-1}\bigr)\bigl(m + 1 - 2w\bigr)
        & (m > v \text{ or } m = v \geq 2w),\ p \text{ split}.
    \end{cases}
  \]
  (Read $\varepsilon(m)\,p^{m/2}$ as $p^{m/2}$ for $m$ even and $0$ for $m$
  odd. For $p \mid f$ one has $v \in \{2w, 2w+1\}$ for $p$ odd and
  $v \in \{2w-2, 2w, 2w+1\}$ for $p = 2$, with $v = 2w-2$ iff $D$ odd; so the
  two regimes of $u_p$ are exhaustive and mutually exclusive.)
  Then for every $n \geq 1$,
  \[
    r_d(n) = \prod_{p \mid n} t_p\bigl(v_p(n)\bigr),
    \qquad
    r_d^{\times}(n) = \prod_{p \mid n} u_p\bigl(v_p(n)\bigr);
  \]
  equivalently the products may run over all primes, since
  $t_p(0) = u_p(0) = 1$.
\end{theorem}

\begin{proof}
  \uses{infra-glue-primary, infra-glue-coprime, prop-3-1-1, thm-3-1-2,
        cor-3-1-3}
  Write $m_p := v_p(n)$. Proposition~\ref{infra-glue-primary} factors both
  counts over the primes dividing $n$:
  $r_d(n) = \prod_{p \mid n} r_d(p^{m_p})$ and
  $r_d^{\times}(n) = \prod_{p \mid n} r_d^{\times}(p^{m_p})$.

  For $p \nmid f$, $\gcd(p^{m_p}, f) = 1$, so
  Theorem~\ref{infra-glue-coprime} gives
  $r_d(p^{m_p}) = r_d^{\times}(p^{m_p}) = r_D(p^{m_p})$, and
  Proposition~\ref{prop-3-1-1} evaluates $r_D(p^m)$ as $m+1$ ($p$ split),
  $\varepsilon(m)$ ($p$ inert), $1$ ($p$ ramified); this is
  $t_p(m_p) = u_p(m_p)$.

  For $p \mid f$, the corrected Theorem~\ref{thm-3-1-2} (statement per
  ERRATA~E4, $p$-uniform) evaluates $r_d(p^{m_p}) = t_p(m_p)$, and the
  corrected Corollary~\ref{cor-3-1-3} (regime per ERRATA~E5) evaluates
  $r_d^{\times}(p^{m_p}) = u_p(m_p)$. The parenthetical trichotomy of $v$
  follows from $d = Df^2$ with $D$ fundamental: for odd $p \mid f$,
  $v = 2w + v_p(D)$ with $v_p(D) \in \{0,1\}$; for $p = 2 \mid f$,
  $v = 2w - 2 + v_2(D)$ with $v_2(D) \in \{0, 2, 3\}$; in particular
  $v = 2w - 1$ never occurs, so at $m = v$ exactly one of $v = 2w-2$,
  $v \geq 2w$ holds.

  Finally, $t_p(0) = u_p(0) = 1$ for every prime $p$: if $p \nmid f$ all three
  cases give $1$ at $m = 0$; if $p \mid f$ then $0 = m \leq v$ gives
  $t_p(0) = (p - 1)/(p - 1) = 1$, and for $u_p$ either $v > 0$ (trivial branch,
  value $\varepsilon(0)p^0 = 1$) or $v = 0$, which forces $p = 2$, $D$ odd,
  $w = 1$, i.e.\ $m = v = 2w - 2$ (again the trivial branch, value $1$).
  Substituting the factor evaluations completes the proof.
\end{proof}

\begin{remark}
  \label{infra-glue-remarks}
  \uses{infra-glue}
  (i) \emph{Divergence from the thesis.} None of the above is stated in the
  thesis, which counts only prime-power norms; and the local factors at
  $p \mid f$ are the ERRATA-corrected statements (E1--E4 for
  Theorem~\ref{thm-3-1-2}, E5 for Corollary~\ref{cor-3-1-3}), not the thesis's
  originals, two of whose $2$-adic branches are provably false. The thesis's
  Corollary~3.1.3 display also reads like a total count but means
  \emph{invertible} ideals (ERRATA E6.1); the notation $r_d$ vs.\
  $r_d^{\times}$ removes the ambiguity.
  (ii) \emph{Index is not multiplicative in $\mathcal{O}_d$.} For $d = -12$,
  $\mathfrak{p} = (2, 1 + \sqrt{-3})$ has index $2$ but $\mathfrak{p}^2$ has
  index $8$; hence Proposition~\ref{infra-glue-primary} is proved via comaximal
  CRT decomposition, never via norm multiplicativity of products (and Mathlib's
  Dedekind-gated \texttt{Ideal.absNorm} API is unusable here).
  (iii) \emph{Coprimality is necessary} in
  Theorem~\ref{infra-glue-coprime}: $r_{-12}(2) = 1$ while $r_{-3}(2) = 0$.
  (iv) \emph{Lean route.} The mod-$N$ comparison of
  Lemma~\ref{infra-glue-avoid} is where Mathlib's conductor machinery applies
  (\texttt{conductor},
  \texttt{Localization.localRingHom\_bijective\_of\_not\_conductor\_le},
  \texttt{comap\_map\_eq\_map\_adjoin\_of\_coprime\_conductor}); a per-prime
  localization proof is a viable alternative to the elementary mod-$N$
  argument given here.
  (v) For $\gcd(n, f) = 1$ the closed form collapses to the classical divisor
  sum $r_d(n) = \sum_{c \mid n} \chi_D(c)$.
  (vi) \emph{Role.} This theorem is the counting instrument for the
  Question~3.0.1 falsification loop: candidate reformulations of the
  Lauter--Viray count $A(N)$ for $\gcd(m, f_1) > 1$ are tested numerically
  against $v_\ell(F(m))$, which requires exact invertible-ideal counts of
  arbitrary composite norm, including norms meeting the conductor.
\end{remark}
